\subsection{An explicit finite five-point input for every retained elliptic coefficient test}

\begin{result}{A.4.39}{Proved}{The two coloured ten-point squares force the exact finite coefficient relation}
Put
\[
 R_0=\langle v,d\mid(vd)^2=1\rangle,
 \qquad S(v)=d,\quad S(d)=v.
                                                               \tag{FEC.1}
\]
For every specified finite quotient $q:R_0\twoheadrightarrow Q$,
the construction below gives an explicitly bounded integer
$m(q)$ and the actual characteristic quotient
\[
 Z(q)=\Gamma_{0,5}/\kappa_{C(q)}(\Gamma_{0,5}),
 \qquad C(q)=(m(q)!)^{(m(q)!)^{45}}.
                                                               \tag{FEC.2}
\]
Write $D=q(d)$ and $V=q(v)$.  The following inclusion
concerns actual gentle elements, with all multipliers retained:
\[
 \begin{aligned}
 &\left\{F=(\lambda,f)\in\widehat{\mathrm{GT}}_0:
                  p(f)=1\text{ in }Z(q)\right\}\\
 &\qquad\subseteq
 \left\{F:
   f(D,V^2)D^{-4c_2(f)}f(V,D^2)=V^{4c_2(f)}\right\}.
 \end{aligned}
                                                               \tag{FEC.3}
\]
Here $p(f)$ is the literal pentagon word (FPA.2), and
$c_2(f)$ is the coefficient defined by the actual index-two
Schreier abelianization in E.22.  Every map defining $Z(q)$
is an ordinary group homomorphism before completion.  Only
the finite pentagon in (FEC.2), the gentle relations, and the
specified finite source image are used.

For every finite homomorphism from the mixed braid group
$\langle u,w^2\rangle\leq B_3$, the same construction, applied
to its exact capped quotient from A.4.28, forces the retained
finite identity
\[
 f(u,w^4)=w^{8c_2(f)}f(u^2,w^2)u^{4c_2(f)}
                         (uw)^{-6c_2(f)}.
                                                               \tag{FEC.4}
\]
The symbols on this last line mean their images in that
specified finite mixed-braid target.  Thus (FEC.2) supplies
the finite test (TA.2), with its central coordinate, at every
prescribed target.  The result does not assert that the
remaining finite shadow census has been executed.
\end{result}

\begin{proof}
\par\medskip\noindent\textbf{A cusp centralizer with a specified finite output.}\quad
Write $a=vd$, so $R_0=\langle a,d\mid a^2=1\rangle$.
Start with a finite quotient $q_0:R_0\to Q_0$ that is
injective on $\langle a\rangle$, and put $K_0=\ker q_0$.
The HNN Bass--Serre tree has vertex stabilizers conjugate to
$\langle a\rangle$, and trivial edge stabilizers.  The action
of $K_0$ is free.  The finite quotient graph $Y=K_0\backslash T$
has vertices $Q_0/\langle q_0(a)\rangle$ and oriented edges
indexed by $Q_0$; its edge indexed by $x$ goes from
$x\langle q_0(a)\rangle$ to
$xq_0(d)\langle q_0(a)\rangle$.
Left deck translation preserves these orientations.
Consequently
\[
 \operatorname{rank}K_0=1+|Q_0|/2,
 \quad Q_1=R_0/(K_0'K_0^2),
 \quad |Q_1|=|Q_0|2^{1+|Q_0|/2}.
                                                               \tag{FEC.5}
\]
These are the actual free-group homology quotient and its
projection to $Q_0$.

Let $n_d=\operatorname{ord}(q_0(d))$.  The element $d^{n_d}$
is a loop in $Y$ whose mod-two edge chain is
\[
          z_d=\sum_{j=0}^{n_d-1}[q_0(d)^j]
                          \quad\in\mathbb F_2[Q_0].
                                                               \tag{FEC.6}
\]
It is nonzero, also when $n_d=1$.  Because a graph has no
two-cells, its first homology injects into its edge-chain
space.  The deck stabilizer of (FEC.6) is exactly
$\langle q_0(d)\rangle$: its support is that subset of
$Q_0$, and its left translates have the same support
precisely for elements of that subgroup.
An element centralizing $d$ in $Q_1$ fixes the image of
$d^{n_d}$ in $K_0/K_0'K_0^2$, so its projection belongs
to that stabilizer.  Conversely every $q_0(d)^j$ is the
projection of the centralizing element $d^j\in Q_1$.
Thus the exact comparison is
\[
 \operatorname{im}\bigl(C_{Q_1}(d)\longrightarrow Q_0\bigr)
                         =\langle q_0(d)\rangle.
                                                               \tag{FEC.7}
\]
Replacing the original quotient by its joint image with
its composite with $S$ makes $K_0$ invariant under $S$.
Then $K_0'K_0^2$ is invariant as well, and applying $S$
to (FEC.7) gives the same assertion for $v$.
Every finite power by $\lambda\in\widehat{\mathbb Z}^{\times}$
generates the same cyclic subgroup, so the centralizers
of $v^\lambda,d^\lambda$ are the same as those of $v,d$.

\par\medskip\noindent\textbf{Two layers determine the coefficient itself.}\quad
Let $b=|Q|$ and $e=\exp(Q)$.  Retain the following actual
finite maps in one joint quotient $q_0:R_0\to Q_0$:
\[
 \begin{gathered}
 q,\quad qS,\quad
 R_0\longrightarrow(\mathbb Z/(2e))[d]\oplus C_2[a],\\
 R_0\longrightarrow R_0/(K_{\rm std}'K_{\rm std}^{e}),\qquad
 R_0\longrightarrow R_0/(K_{\rm geo}'K_{\rm geo}^{e}),\\
 K_{\rm std}=F(v,d^2),\qquad K_{\rm geo}=S(K_{\rm std})=F(d,v^2).
 \end{gathered}
                                                               \tag{FEC.8}
\]
The free-group identities for these index-two kernels
are the actual Reidemeister--Schreier identities (A.T21).
Their displayed characteristic subgroups are normal in
$R_0$, so the last two lines are finite group quotients,
each of order $2e^2$.  The abelian quotient has order
$4e$ and is $S$-invariant: on it $S(d)=a-d$ and $S(a)=a$.
The first two coordinates and last two coordinates are
interchanged by $S$.  Thus $\ker q_0$ is $S$-invariant,
$q_0$ retains the cone involution, and
\[
                         |Q_0|\leq16b^2e^5.
                                                               \tag{FEC.9}
\]
Starting from this $Q_0$, form the two homology refinements
\[
 Q_1=R_0/(K_0'K_0^2),\quad K_1=\ker(R_0\to Q_1),
 \qquad Q_2=R_0/(K_1'K_1^2).
                                                               \tag{FEC.10}
\]
Both projections have the cusp-centralizer property
(FEC.7), for both cusps.

Put $h=f(v,d^2)$ and $k=f(d,v^2)$, always evaluated
through the named finite homomorphisms.  Consider the
two actual gentle mixed actions
\[
 \begin{array}{ll}
 \phi(v)=v^\lambda,&\phi(d)=h^{-1}d^\lambda h,\\
 \psi=S\phi S,&
 \psi(d)=d^\lambda,\quad\psi(v)=k^{-1}v^\lambda k.
 \end{array}
                                                               \tag{FEC.11}
\]
The geometric construction below gives their inner
equivalence at a universal finite quotient mapping to
$Q_2$.  Project that equality to $Q_2$ and write its
conjugator as $g$ there.  The equation on $d$ says
$gh^{-1}\in C_{Q_2}(d^\lambda)$.  The first use of
(FEC.7), down to $Q_1$, gives an integer $t$ with
$g=d^th$ at $Q_1$.  The equation on $v$ now says
$kd^th\in C_{Q_1}(v^\lambda)$.  The second use of
(FEC.7), down to $Q_0$, gives an integer $s$ with
\[
                         kd^th=v^s\quad\text{in }Q_0.
                                                               \tag{FEC.12}
\]
No asserted centralizer equality in $Q_0$ is used.

Both $h$ and $k$ have zero $R_0$ abelianization.
The retained third coordinate of (FEC.8) therefore gives
$t=-s\pmod{2e}$ and $s=0\pmod2$.  In particular the
integer $t$ is even.  The whole equality (FEC.12) is
then in the image of $K_{\rm std}$, where its retained
abelian coordinates are computed in the basis
$X=v$, $Y=d^2$.
Write $c=c_2(f)$.  In the actual index-two subgroup of
$\widehat F(x,y)$ with $x$ odd and $y$ even, the exponent
vector of $f$ on the basis $(x^2,y,xyx^{-1})$ is
$(0,c,-c)$.  Under $(x,y)\mapsto(d,v^2)$ this basis
maps to
\[
                    Y,\quad X^2,\quad(X^{-1}Y^{-1})^2.
                                                               \tag{FEC.13}
\]
The last equality uses $dvd^{-1}=v^{-1}d^{-2}$ in
the displayed presentation of $R_0$.
Consequently $[k]=(4c,2c)$ and $[h]=0$ in
$\widehat{K_{\rm std}^{\rm ab}}$.  Applying the retained mod-$e$
quotient to (FEC.12) gives
\[
           (4c,2c+t/2)=(s,0)\pmod e,
 \qquad t=-4c\pmod{2e},\quad s=4c\pmod e.
                                                               \tag{FEC.14}
\]
Projection to the original exponent-$e$ target $Q$
is exactly the equation on the right of (FEC.3).
The exponent $t$ has been determined by an actual
retained quotient, rather than by a formal division
inside a finite cyclic group.

\par\medskip\noindent\textbf{The genuine mixed action at every multiplier.}\quad
For completeness, the representatives (FEC.11) exist
for every gentle element before any finite pentagon
is imposed.  Use its genuine modular action on
$B=(C_2*C_3)^{\wedge}$, with
$\tau=sr$, $x=\tau^{-2}$, $y=sxs$,
$\alpha(\tau)=\tau^\lambda$, and
$\alpha(y)=f(x,y)^{-1}y^\lambda f(x,y)$.
The actual index-three subgroup
$R_H=\langle x,s\rangle$ has the marked isomorphism
\[
 j:R\xrightarrow{\sim}R_H,
 \quad j(v)=x,\quad j(d)=x^{-1}s,
 \quad j(d^2)=x^{-1}y^{-1}.
                                                               \tag{FEC.15}
\]
Here $R=\widehat R_0$.
Put $m_F=(\lambda-1)/2$ and
$\theta=\operatorname{Inn}(\tau)|_{\widehat F(x,y)}$.
The literal free-product identities give
$\theta(x)=x$, $\theta(y)=x^{-1}y^{-1}$ and
\[
 \alpha\theta=\operatorname{Inn}(x^{-m_F})\theta\alpha,
 \qquad
 j\phi j^{-1}=\operatorname{Inn}(x^{m_F})\alpha|_{R_H}.
                                                               \tag{FEC.16}
\]
To check the second identity, the first computes both
images on the free basis $(x,\theta(y))$ and therefore
the images of $v,d^2$.  The two proposed images of $d$
have the same square.  After conjugation by $h$ they
centralize $d^{2\lambda}$.  In $R$ that centralizer is
$\langle d\rangle$: after writing an element as $d^\epsilon k$
with $k\in\widehat K_{\rm std}$, its $k$ part centralizes the
primitive generator $d^2$ of that free profinite group.
The primitive peripheral centralizer is procyclic, as
proved by the finite-cover argument already used in A.1.
Multiplication by two is injective on $\widehat{\mathbb Z}$,
so equality of the two squares gives equality of the
images of $d$.  This proves (FEC.16) and (FEC.11).
Finally $m_{FF'}=m_F+\lambda_Fm_{F'}$ shows that this
choice is an actual action in $\operatorname{Aut}(R)$.
Its only input is the genuine gentle modular action.

The finite four-point action supplied by A.4.37 has
this same outer representative, including its deck
quotient.  This comparison can be checked on all
three sphere braid generators.  The actual projection
\[
 \Gamma_{0,[4]}\longrightarrow B_3/Z(B_3)=B_0
 \quad\text{sends}\quad
       (\sigma_1,\sigma_2,\sigma_3)\longmapsto(u,w,u).
                                                               \tag{FEC.17}
\]
Here $B_0=C_2*C_3$ is discrete and $B=\widehat B_0$.
Its kernel is the ordinary Klein four group of double
transpositions.  In its target,
$y_3=w u^2 w=u^{-2}$, since
$u^2y_3=(uw)^3=1$.  Thus the third standard braid
conjugator becomes $f(u^2,u^{-2})=1$.  The three
images of the finite action project to
$u^\lambda$, $f(u^2,w^2)^{-1}w^\lambda f(u^2,w^2)$,
$u^\lambda$, respectively.  These are exactly the
modular action on the indicated three generators.
This proves equivariance of the actual quotient
(FEC.17) by substitution; no centralizer uniqueness
at a finite level is inferred from agreement on its
pure subgroup.  There is also a direct check in the coloured braid marking
$d=u$, $v=w^2$ from (A.J13).  Its raw restriction is
$d\mapsto d^\lambda$, $v\mapsto h v^\lambda h^{-1}$,
where $h=f(v,d^2)$, by antisymmetry.  Conjugation by
$h^{-1}$ gives exactly $\phi$ in (FEC.11).  Thus the
coloured marking is checked before taking a finite
quotient; (FEC.16) gives its separate comparison with
the modular $R_H$ marking.

\par\medskip\noindent\textbf{The finite square, with every intermediate map retained.}\quad
Let $T(n)=n2^{1+n/2}$ for even positive $n$.  Define
\[
 s_0=16b^2e^5,\quad M=T(T(s_0)),\quad
 R_* =R_0/\kappa_M(R_0),\quad
 B_*=(M!)^{(M!)^2}.
                                                               \tag{FEC.18}
\]
Equations (FEC.5), (FEC.9) and (FEC.10) give a literal
projection $R_*\twoheadrightarrow Q_2$, because the
regular representation of $Q_2$, padded by fixed
letters to degree $M$, is one of the defining maps.
Also $|R_*|\leq B_*$, since $R_0$ has two generators.
The kernel defining $R_*$ is invariant under all
automorphisms of $\widehat R_0$, in particular $\phi$
and $S$.  Thus both in (FEC.11) are automorphisms
of this same finite group.

Use exactly the ten-point markings (A.T10), the
group $G=G_3$ of (HCF.2), its $H$, and
$N=C_G(H)$.  Let $M_{10}=\Gamma_{0,[10]}$.
There are two specified finite-index subgroups of
$M_{10}$: $M_b$ preserves the two labelled blocks
$A_0$ and $O_b=\{z,1/z\}$ separately, and $M_a$
preserves $B_0$ and $O_a=\{z,-z\}$ separately.
Each has index
$\binom{10}{2}\binom82=1260$.
Forgetting to these respective quartets, then using
the actual coloured quotient (A.T16), and finally
projecting to $R_*$ gives the homomorphisms
\[
       \rho_b:M_b\longrightarrow R_* ,\qquad
       \rho_a:M_a\longrightarrow R_* .
                                                               \tag{FEC.19}
\]
Induce each map to $M_{10}$ on its actual finite
coset set.  Its value group is a subgroup of
$R_*^{1260}\rtimes\mathfrak S_{1260}$, so has order
at most $1260! B_*^{1260}$.  At the identity coset,
the induced map restricted to $M_i$ recovers precisely
$\rho_i$.  Put
\[
 m_0=1260! B_*^{1260},\qquad
 L_0=S_{10}(m_0),\qquad \eta_0:M_{10}\to L_0.
                                                               \tag{FEC.20}
\]
Both induced finite maps factor through $\eta_0$ by
their regular representations.  In particular each
$\rho_i$ is an actual map on the specified subgroup
$\eta_0(M_i)$, not an unidentified quotient of $N$.

Apply A.4.34 to $\eta_0|_G$, giving the centralizer
refinement $G\to C$.  Apply A.4.35 to that quotient,
giving $G\to D$.  These maps have the specified
quotient projections
\[
                 D\longrightarrow C\longrightarrow \eta_0(G).
                                                               \tag{FEC.21}
\]
Induce the $D$ map to $M_{10}$ on its
$t_G=10!/32$ actual cosets of $G$.  Let $d_*$ be
the explicit upper bound for $|D|$ obtained by the
three-stage recursions (HCF.11)--(HCF.12) and
(KFU.11), starting with
\[
 |\eta_0(G)|\leq32(m_0!)^{(m_0!)^{45}}.
\]
Define, without any search for an unknown level,
\[
                  m(q)=\max\{m_0,t_G!d_*^{t_G}\}.
                                                               \tag{FEC.22}
\]
This is a finite expression in the specified output
order and the displayed three-step integer recursions.
The universal input $L_1=S_{10}(m(q))$ maps both
to $L_0$ and to the induced quotient, whose identity
coordinate retains (FEC.21).

For an element of the locus on the left of (FEC.3),
A.4.37 constructs compatible automorphisms $\Phi_1$
and $\Phi_0$ on $L_1$ and $L_0$, with the ordinary
labelled forgetting maps.  They induce the identity
on label permutations.  Project $\Phi_1(H)$ to $D$.
It is a homomorphism from the named $H$ with its
named elementwise labels.  A.4.35 supplies an actual
$g\in G$ which conjugates its image in $C$ pointwise
to the actual image of $H$.  Choose the corrector in
the form given there: a pure mapping class times at
most the half twist supported on $C_0$.
Replace both automorphisms by
$\Phi'_j=\operatorname{Inn}(g)\Phi_j$.
The retained $\mathfrak S_{10}$ quotient and the fact
that $G$ is its full inverse image of
$C_{\mathfrak S_{10}}(H)$ imply that $\Phi'_1$
preserves the image of $G$ in $L_1$.  Its projections
to $D$ and $C$ are therefore defined on every element
used here.

For every ordinary $n\in N$, its $\Phi'_1$ image
commutes with the images of $H$.  At $C$ these are
the named elements of $H$.  A.4.34 therefore places
its lower image $\Phi'_0(\eta_0(n))$ inside $\eta_0(N)$.
Since $\Phi'_0$ is an automorphism of the finite
group $L_0$, this inclusion is equality on the
finite subgroup $\eta_0(N)$.  Its action fixes $\eta_0(H)$
pointwise, and hence induces an automorphism
\[
       \alpha:L\xrightarrow{\sim}L,
       \qquad L=\eta_0(N)/\eta_0(H).
                                                               \tag{FEC.23}
\]
This is where invertibility and the actual geometric
centralizer image are both used.  No invariance of
the kernels of $C$ or $D$ under $\Phi_1$ is needed.

The maps (FEC.19) induce surjections
$\overline q_b,\overline q_a:L\to R_*$, with
\[
                       \overline q_a=S\overline q_b.
                                                               \tag{FEC.24}
\]
Here is the precise descent.  By (A.T17),
$N=HN_i$ and $N_i\cap H$ is the named order-two
double flip killed by $\rho_i$.
The retained label map implies
$\eta_0(H)\cap \eta_0(N_i)=\eta_0(H\cap N_i)$.
Thus one may represent each class in $L$ by an
element of $\eta_0(N_i)$, and apply the identity
coordinate of the induced map (FEC.19).  Its kernel
contains $\ker(\eta_0|_{N_i})$ by (FEC.20), and it
kills the indicated intersection with $\eta_0(H)$.
This proves the claimed descent.
Identity (FEC.24) is the projection of the literal
ordinary equality (A.T19), with its chosen
basepoint comparison.  Surjectivity follows from
the ordinary map (A.T20).

The corrector $g$ belongs to both $M_a$ and $M_b$.
Its possible $C_0$ half twist is killed by both
forgetting maps; its pure part contributes the
actual target inner automorphisms.
Finite forgetting compatibility from A.4.37,
the explicit four-point projection (FEC.17), and
the common marked representative (FEC.16) therefore
give elements $g_a,g_b\in R_*$ with
\[
 \overline q_b\alpha
       =\operatorname{Inn}(g_b)\phi\overline q_b,
 \qquad
 \overline q_a\alpha
       =\operatorname{Inn}(g_a)\phi\overline q_a.
                                                               \tag{FEC.25}
\]
The retained labels and the pure/$C_0$-swap
permutation of $g$ show that $\Phi'_0$ preserves
$\eta_0(N_i)$: it preserves $\eta_0(N)$, and $N_i$
is its subgroup specified by the retained orbit
stabilizer.  Thus (FEC.25) follows first on each
$N_i$ from its actual forgotten quartet, and then
on $L$ from the descent just proved.  The two inner elements are retained
separately.  Substituting (FEC.24) in (FEC.25) and
cancelling the surjection $\overline q_b$ yields
\[
       S\phi=\operatorname{Inn}
          \bigl(S(g_b)^{-1}g_a\bigr)\phi S
                      \quad\text{in }\operatorname{Aut}(R_*).
                                                               \tag{FEC.26}
\]
Composing on the right by $S$ gives precisely the
inner equivalence of $\psi=S\phi S$ and $\phi$
used before (FEC.12).  Projecting through the actual
map $R_*\to Q_2\to Q_1\to Q_0\to Q$ and using
(FEC.12)--(FEC.14) proves (FEC.3).

\par\medskip\noindent\textbf{Retaining the braid center.}\quad
Let $q_H$ be the specified finite mixed-braid map
on $H=\langle u,w^2\rangle$, and let $z=(uw)^3$.
Its abstract presentation is
$H=\langle d,V\mid(dV)^2=(Vd)^2\rangle$ with
$d=u$, $V=w^2$, and $z=(dV)^2$.
Apply the exact construction (MCR.2): join $q_H$
with the character $d\mapsto1,V\mapsto0$ modulo
$2\operatorname{ord}(q_Hz)$ and kill the diagonal
image of $z$.  This gives an actual quotient of
$R_0$ of order at most $2|q_H(H)|$, with $v=V$.
Use that quotient as the input $q$ in (FEC.2).

Antisymmetry gives $f(d^2,v)=f(v,d^2)^{-1}$.
Therefore the equation on the right of (FEC.3)
is exactly
$f(d,v^2)=v^{4c}f(d^2,v)d^{4c}$ in the capped
target.  A.4.28 recovers the central factor
$z^{-2c}$ in the original finite mixed target,
because the whole defect has zero character.
Substituting $d=u,V=w^2,z=(uw)^3$ gives (FEC.4).
This verifies the actual morphism between the finite
orbifold and braid equations, including the exponent
of their common central generator.

The independent program
\path{work/tower_code/finite_cusp_cycle.py}
checks both cusp cycles for all $216$ homomorphisms to
$\mathfrak S_4$ retaining a nontrivial cone involution.
It records the literal graph vertices, oriented edges,
cycle supports, all deck stabilizers, homology ranks
and exact extension orders.  It also checks the
free-product words (FEC.13) and (FEC.17).  These finite
examples verify the displayed maps; the general
centralizer assertion is the graph proof (FEC.5)--(FEC.7).
The executed source, output and JSON data are retained.

All quotients and all induced maps in the proof are
given by finite group presentations, homology of
specified finite graph covers, or intersections of
the kernels of homomorphisms to a specified finite
symmetric group.  The argument gives a terminating
finite construction with the displayed bound.
It proves the inclusion (FEC.3) for the actual gentle
source image and does not replace that image by all
formal pairs satisfying a collection of finite
residues.  Computing the complete source image and
the final pentagon shadow census remains the next
step required by A.4.
\end{proof}
